\begin{center}
\vspace*{1\baselineskip}
{\fontsize{14pt}{17pt}\selectfont\textbf{Abstract}}
\vspace{1\baselineskip}
\end{center}

\begin{singlespace}
\noindent
This thesis presents the design, implementation, and empirical evaluation of an optimized modular computational architecture for high-throughput engineering workflows.
Modern distributed environments often suffer from structural friction, version collisions, and synchronization bottlenecks during collaborative research and deployment cycles.
To address these critical challenges, our proposed methodology decouples procedural system components, establishes isolated mathematical formulations, and enforces reproducible data pipelines across heterogeneous runtime nodes.
Experimental investigations conducted across benchmark workloads demonstrate that the proposed architecture yields statistically significant improvements in processing latency and operational precision compared to baseline implementations.
Specifically, our framework achieves a 95.7\% precision and a 95.2\% overall F1-score while preserving strict formatting determinism and low resource overhead.
These findings validate the practical viability of the decoupled architecture for real-world engineering and computational tasks.

\vspace{2\baselineskip}

\noindent\textbf{Keywords (maximum of five):} modular architecture, empirical evaluation, computational optimization, distributed systems, reproducible pipelines
\end{singlespace}
\newpage
